% J16 — A Commutative Non-Associative 4-Algebra over F_5 with Rigid
%       Idempotent Decomposition (was: "Discrete Dirac on F_5^4")
%
% Brayden R. Sanders, M. Gish, 7Site LLC
% Submitted to: Algebras and Representation Theory
% Reference impl: Gen13/targets/ck/brain/dirac/tig_dirac.py
% Verification:   verify_discrete_dirac_4core.py + test_tig_dirac.py
% MSC 2020: 17A30, 15A66, 17A40, 11T55

\documentclass[11pt,reqno]{amsart}

\usepackage{amsmath, amssymb, amsthm, mathtools}
\usepackage[margin=1in]{geometry}
\usepackage[colorlinks=true,linkcolor=blue,citecolor=blue,urlcolor=blue]{hyperref}
\usepackage{microtype}
\usepackage{booktabs}
\usepackage{enumitem}

\theoremstyle{plain}
\newtheorem{theorem}{Theorem}[section]
\newtheorem{proposition}[theorem]{Proposition}
\newtheorem{lemma}[theorem]{Lemma}
\newtheorem{corollary}[theorem]{Corollary}

\theoremstyle{definition}
\newtheorem{definition}[theorem]{Definition}

\theoremstyle{remark}
\newtheorem*{remark}{Remark}

\newcommand{\Z}{\mathbb{Z}}
\newcommand{\F}{\mathbb{F}}
\newcommand{\R}{\mathbb{R}}
\newcommand{\Cl}{\mathrm{Cl}}
\newcommand{\Aut}{\mathrm{Aut}}

\title[A 4-algebra over $\F_5$ with rigid idempotent decomposition]
{A Commutative Non-Associative 4-Algebra over $\F_5$\\
 with Rigid Idempotent Decomposition}

\author{Brayden R.\ Sanders}
\address{7Site LLC, Hot Springs, Arkansas, USA}
\email{brayden@7site.co}

\author{M.\ Gish}
\address{Independent Researcher, Hot Springs, Arkansas, USA}

\subjclass[2020]{17A30, 15A66, 17A40, 11T55}
\keywords{commutative non-associative algebra, finite-field algebra,
idempotent decomposition, automorphism group, rigid algebraic structure}

\date{\today}

\begin{document}

\begin{abstract}
We study a four-dimensional commutative non-associative algebra $V$ over the prime field $\F_5$, defined by the bilinear extension of an explicit $4 \times 4$ multiplication table on a basis $\{e_0, e_2, e_3, e_4\}$. The algebra has a remarkably rigid structure. Its idempotent set has exactly four elements: $0$, two non-zero \emph{primitive} idempotents $p_+ = e_2$ and $p_- = e_0 + 4 e_2 = e_0 - e_2 \pmod 5$ satisfying $p_+ p_- = 0$, and the third non-zero idempotent $e_0 = p_+ + p_-$. The algebra has a one-dimensional left annihilator ideal generated by $\varepsilon = e_3 - e_4 \pmod 5$. The left-multiplication operators $L_{e_2}$ and $L_{e_0}$, viewed as $\F_5$-linear maps $V \to V$, have eigenspace-dimension profiles $(1, 3)$ and $(2, 2)$ respectively, the simultaneous eigenspace $\bigl(\text{1-eigenspace of } L_{e_2}\bigr) \cap \bigl(\text{0-eigenspace of } L_{e_0}\bigr)$ is trivial, and the two operators commute. The associator $[x, y, z] := (xy)z - x(yz)$ has image contained in the one-dimensional subspace $\F_5 \cdot p_-$. The algebra is power-associative. The automorphism group has order $40$, and we identify its structure as a non-trivial extension involving the Frobenius group $F_{20}$ and a $\Z/2$ factor. All structural claims are verified independently by short Python scripts ($< 2$ sec total runtime, \texttt{numpy} only) which we describe in §\ref{sec:verification}.
\end{abstract}

\maketitle

\tableofcontents

%==============================================================
\section{Introduction}
\label{sec:intro}
%==============================================================

\subsection{Scope and tier discipline}

\textbf{Proven:} The structural theorem (Theorem~\ref{thm:struct}) characterizes the four idempotents of $V$, the eigenspace-dimension profiles of $L_{e_2}$ and $L_{e_0}$, the commutation $L_{e_2} L_{e_0} = L_{e_0} L_{e_2}$, the trivial simultaneous eigenspace, the one-dimensional associator image, power-associativity, and the order-$40$ automorphism group with structure described in Theorem~\ref{thm:aut}. \textbf{Computed:} All structural claims independently verified by 14 algebraic checks in \texttt{verify\_discrete\_dirac\_4core.py} and 15 unit tests in \texttt{test\_tig\_dirac.py}; total runtime under 2 seconds with \texttt{numpy} as the only dependency. \textbf{Structural rhyme:} The eigenspace profiles $(1, 3)$ and $(2, 2)$ and the empty simultaneous eigenspace are reminiscent of the $1+3$ Minkowski signature and $2+2$ chirality split of the relativistic Dirac equation, with the empty simultaneous eigenspace echoing the V--A asymmetry of the weak interaction. These are interpretive analogies; we make no claim of physical interpretation. The algebra is in the broader theme of rigid-idempotent commutative algebras (e.g., axial algebras of Hall--Rehren--Shpectorov~\cite{HallRehrenShpectorov2015}), though $V$ is not technically an axial algebra (it lacks the standard Frobenius form). \textbf{Open:} Field-invariance of the structural features over arbitrary primes $p \notin \{2, 5\}$; explicit characterization of higher tensor powers; whether $V$ admits an action of $\mathrm{SU}(5)$ or related compact groups in any compatible sense. The framing follows the Drápal--Wanless~\cite{DrapalWanless2021} line of work on small finite commutative non-associative structures.

\subsection{Lens-ownership}

This paper works over the prime field $\F_5$ with a specific 4-dimensional commutative non-associative algebra $V$ defined by the bilinear extension of an explicit multiplication table on a designated basis (Definition~\ref{def:V}). The choice of this specific table is foundational, not derived: it arose historically as the bilinear lift to $\F_5$ of a 4-element subset of $\Z/10\Z$ under a related composition (cited from a companion submission~\cite{SandersGishFourCore}), but the math of this paper does not depend on that origin. The theorems below are theorems on this specific algebra; analogous theorems would hold for other 4-dimensional commutative non-associative $\F_p$-algebras with similar rigid idempotent structure.

\subsection{Main result}

The algebra is defined by Definition~\ref{def:V} below. The single main result is:

\begin{theorem}[Structural rigidity of $V$, abridged statement] \label{thm:struct-abridged}
Let $V$ be the algebra of Definition~\ref{def:V}. Then:
\begin{enumerate}[label=\textup{(\alph*)}]
\item $V$ has exactly four idempotents: $0$, $p_+ := e_2$, $p_- := e_0 + 4 e_2 \pmod 5$, and $e_0 = p_+ + p_-$.
\item $V$ has a one-dimensional left annihilator ideal, namely $\F_5 \cdot \varepsilon$ where $\varepsilon := e_3 - e_4 \pmod 5$.
\item $L_{e_2}$ has eigenspace-dimension profile $(1, 3)$ on $V$ (its $1$- and $0$-eigenspaces have $\F_5$-dimensions $1$ and $3$ respectively); $L_{e_0}$ has profile $(2, 2)$.
\item The simultaneous eigenspace $\bigl(\text{1-eigen of } L_{e_2}\bigr) \cap \bigl(\text{0-eigen of } L_{e_0}\bigr)$ is trivial.
\item $L_{e_2}$ and $L_{e_0}$ commute as $\F_5$-linear maps.
\item For all $x, y, z \in V$, the associator $[x, y, z]$ lies in $\F_5 \cdot p_-$.
\item $V$ is power-associative: $(xx) x = x (xx)$ for all $x \in V$.
\item No $\F_5$-algebra automorphism of $V$ swaps $p_+$ and $p_-$.
\item $|\Aut(V)| = 40$, with structure $\Aut(V) \cong F_{20} \rtimes \Z/2$ (more precisely, a group of order $40$ containing $F_{20}$ as a normal subgroup of index $2$).
\end{enumerate}
\end{theorem}

The detailed proof is in §\ref{sec:proofs}.

\subsection{Related companion work}

A companion submission~\cite{SandersGishFpUniversality} extends the structural features (a)--(g) above from $\F_5$ to $\F_p$ for $p \in \{2, 3, 5, 7, 11, 13\}$ via the same bilinear-extension construction with table entries reduced mod $p$. A second companion~\cite{SandersGishCliffordLadder} establishes a dimensional ladder $\dim_{\F_5} V^{\otimes n} = 4^n = \dim_\R \Cl(2n)$ for $n \leq 5$ together with a refined-cell binomial grading matching $\Cl(2n)$ grade dimensions. These are stated only for context; the present paper is self-contained.

%==============================================================
\section{The algebra $V$}
\label{sec:V}
%==============================================================

\begin{definition}\label{def:V}
Let $V$ be the four-dimensional $\F_5$-vector space with basis $\{e_0, e_2, e_3, e_4\}$. The indexing reflects the historical origin (a 4-element subset $\{0, 7, 8, 9\} \subset \Z/10\Z$ projected mod 5), but the mathematics in this paper depends only on the explicit table below. Define multiplication $\cdot: V \times V \to V$ as the bilinear extension over $\F_5$ of the table
\begin{equation}\label{eq:mult-table}
\begin{array}{c|cccc}
\cdot & e_0 & e_2 & e_3 & e_4 \\ \hline
e_0 & e_0 & e_2 & e_0 & e_0 \\
e_2 & e_2 & e_2 & e_2 & e_2 \\
e_3 & e_0 & e_2 & e_2 & e_2 \\
e_4 & e_0 & e_2 & e_2 & e_2
\end{array}
\end{equation}
where the $(i, j)$ entry is $e_i \cdot e_j$. We refer to the algebra $(V, \cdot)$ simply as $V$.
\end{definition}

\begin{remark}[Commutativity]
A direct check of equation~\eqref{eq:mult-table} shows the table is symmetric under $(i, j) \leftrightarrow (j, i)$ (using the identifications $e_0 \cdot e_2 = e_2 = e_2 \cdot e_0$, $e_0 \cdot e_3 = e_0 = e_3 \cdot e_0$, $e_0 \cdot e_4 = e_0 = e_4 \cdot e_0$, $e_2 \cdot e_3 = e_2 = e_3 \cdot e_2$, $e_2 \cdot e_4 = e_2 = e_4 \cdot e_2$, $e_3 \cdot e_4 = e_2 = e_4 \cdot e_3$). Hence the bilinear extension is commutative.
\end{remark}

\subsection{Privileged elements}

In addition to the basis $\{e_0, e_2, e_3, e_4\}$, several derived elements play a central role:
\begin{align*}
p_+ &:= e_2 && \text{(``primary primitive idempotent'')},\\
p_- &:= e_0 + 4 e_2 \pmod 5 = e_0 - e_2 \pmod 5 && \text{(``secondary primitive idempotent'')},\\
\varepsilon &:= e_3 + 4 e_4 \pmod 5 = e_3 - e_4 \pmod 5 && \text{(``annihilator nilpotent'')},\\
h &:= 2(e_3 + e_4) \pmod 5 && \text{(``square root of } p_+ \text{'')}.
\end{align*}

These are not independently axiomatic; they are particular elements of $V$, and the structural theorem will recover their roles.

%==============================================================
\section{Proof of the structural theorem}
\label{sec:proofs}
%==============================================================

We prove Theorem~\ref{thm:struct-abridged} directly from the multiplication table~\eqref{eq:mult-table}.

\subsection{The left-multiplication matrices}

For any $a \in V$, $L_a \colon V \to V$ is the $\F_5$-linear map $L_a(v) := a \cdot v$. Writing $a = \sum_i \alpha_i e_i$, $L_a$ is the linear extension of $L_a(e_j) = \sum_i \alpha_i (e_i \cdot e_j)$.

\begin{lemma} \label{lem:L-matrices}
Relative to the ordered basis $(e_0, e_2, e_3, e_4)$, the left-multiplication matrices for $a = e_0$ and $a = e_2$ are
\begin{equation}\label{eq:L-e0-e2}
L_{e_0} = \begin{pmatrix}
1 & 0 & 1 & 1 \\
0 & 1 & 0 & 0 \\
0 & 0 & 0 & 0 \\
0 & 0 & 0 & 0
\end{pmatrix}, \qquad
L_{e_2} = \begin{pmatrix}
0 & 0 & 0 & 0 \\
1 & 1 & 1 & 1 \\
0 & 0 & 0 & 0 \\
0 & 0 & 0 & 0
\end{pmatrix}.
\end{equation}
\end{lemma}

\begin{proof}
The columns of $L_a$ are $L_a(e_0), L_a(e_2), L_a(e_3), L_a(e_4)$ expressed in the basis. Reading off the table~\eqref{eq:mult-table}: $L_{e_0}(e_0) = e_0$, $L_{e_0}(e_2) = e_2$, $L_{e_0}(e_3) = e_0$, $L_{e_0}(e_4) = e_0$, giving columns $(1,0,0,0), (0,1,0,0), (1,0,0,0), (1,0,0,0)$. Similarly $L_{e_2}(e_0) = e_2$, $L_{e_2}(e_2) = e_2$, $L_{e_2}(e_3) = e_2$, $L_{e_2}(e_4) = e_2$, giving columns $(0,1,0,0)$ each.
\end{proof}

\subsection{Eigenspace dimensions}

\begin{lemma} \label{lem:eigenspaces}
\begin{enumerate}[label=\textup{(\alph*)}]
\item $L_{e_2}$ has $\F_5$-rank $1$. The $1$-eigenspace of $L_{e_2}$ has dimension $1$ (spanned by $e_2$), and the $0$-eigenspace has dimension $3$.
\item $L_{e_0}$ has $\F_5$-rank $2$. The $1$-eigenspace of $L_{e_0}$ has dimension $2$, and the $0$-eigenspace has dimension $2$.
\item The simultaneous eigenspace $\bigl(\text{1-eigen of } L_{e_2}\bigr) \cap \bigl(\text{0-eigen of } L_{e_0}\bigr)$ is the zero subspace.
\end{enumerate}
\end{lemma}

\begin{proof}
(a) From~\eqref{eq:L-e0-e2}, $L_{e_2}$ has all rows zero except row $2$ which is $(1, 1, 1, 1)$, so rank $1$. The $1$-eigenspace is the kernel of $L_{e_2} - I$, where
\[
L_{e_2} - I = \begin{pmatrix}
-1 & 0 & 0 & 0 \\
1 & 0 & 1 & 1 \\
0 & 0 & -1 & 0 \\
0 & 0 & 0 & -1
\end{pmatrix} = \begin{pmatrix}
4 & 0 & 0 & 0 \\
1 & 0 & 1 & 1 \\
0 & 0 & 4 & 0 \\
0 & 0 & 0 & 4
\end{pmatrix} \pmod 5.
\]
This matrix has rank $3$ (rows $1$, $3$, $4$ pivot at columns $1$, $3$, $4$; row $2$ is a linear combination), so its kernel has dimension $1$, spanned by $(0, 1, 0, 0)^T = e_2$. The $0$-eigenspace is the kernel of $L_{e_2}$, the set of $v$ with $v_1 = -(v_0 + v_2 + v_3) \pmod 5$, of dimension $3$ (parametrized by $v_0, v_2, v_3$ free).

(b) From~\eqref{eq:L-e0-e2}, $L_{e_0}$ has rows $1, 2$ non-trivial and rows $3, 4$ zero. Row $1$ is $(1, 0, 1, 1)$, row $2$ is $(0, 1, 0, 0)$. These are linearly independent, so rank $2$. The $0$-eigenspace (kernel of $L_{e_0}$) has dimension $4 - 2 = 2$. The $1$-eigenspace is the kernel of
\[
L_{e_0} - I = \begin{pmatrix}
0 & 0 & 1 & 1 \\
0 & 0 & 0 & 0 \\
0 & 0 & -1 & 0 \\
0 & 0 & 0 & -1
\end{pmatrix} \pmod 5.
\]
This has rank $2$ (rows $3$ and $4$ pivot, row $1$ is their sum negated), so kernel dimension $2$.

(c) The $1$-eigenspace of $L_{e_2}$ is $\F_5 \cdot e_2$. For $v = \alpha e_2$ to be in the $0$-eigenspace of $L_{e_0}$, we need $L_{e_0}(\alpha e_2) = \alpha e_2 = 0$, hence $\alpha = 0$. So the simultaneous eigenspace is $\{0\}$.
\end{proof}

\subsection{Commutation of $L_{e_0}$ and $L_{e_2}$}

\begin{lemma} \label{lem:commute}
$L_{e_0} L_{e_2} = L_{e_2} L_{e_0}$ as $\F_5$-linear maps on $V$.
\end{lemma}

\begin{proof}
Direct computation from~\eqref{eq:L-e0-e2}:
\[
L_{e_0} L_{e_2} = \begin{pmatrix}
1 & 0 & 1 & 1 \\
0 & 1 & 0 & 0 \\
0 & 0 & 0 & 0 \\
0 & 0 & 0 & 0
\end{pmatrix}
\begin{pmatrix}
0 & 0 & 0 & 0 \\
1 & 1 & 1 & 1 \\
0 & 0 & 0 & 0 \\
0 & 0 & 0 & 0
\end{pmatrix} = \begin{pmatrix}
0 & 0 & 0 & 0 \\
1 & 1 & 1 & 1 \\
0 & 0 & 0 & 0 \\
0 & 0 & 0 & 0
\end{pmatrix},
\]
and
\[
L_{e_2} L_{e_0} = \begin{pmatrix}
0 & 0 & 0 & 0 \\
1 & 1 & 1 & 1 \\
0 & 0 & 0 & 0 \\
0 & 0 & 0 & 0
\end{pmatrix}
\begin{pmatrix}
1 & 0 & 1 & 1 \\
0 & 1 & 0 & 0 \\
0 & 0 & 0 & 0 \\
0 & 0 & 0 & 0
\end{pmatrix} = \begin{pmatrix}
0 & 0 & 0 & 0 \\
1 & 1 & 1 & 1 \\
0 & 0 & 0 & 0 \\
0 & 0 & 0 & 0
\end{pmatrix}.
\]
\end{proof}

\subsection{Idempotents}

\begin{lemma} \label{lem:idempotents}
The idempotents of $V$ are exactly $\{0, e_0, e_2, e_0 + 4 e_2\}$ (where $e_0 + 4 e_2 = e_0 - e_2 \pmod 5$). The non-zero ones are $e_0, e_2, p_-$ where $p_- := e_0 - e_2$.
\end{lemma}

\begin{proof}
Write $x = \alpha e_0 + \beta e_2 + \gamma e_3 + \delta e_4$. Compute $x^2 = x \cdot x$ using bilinearity and~\eqref{eq:mult-table}:
\begin{align*}
x^2 &= \alpha^2 (e_0 \cdot e_0) + \alpha \beta (e_0 \cdot e_2 + e_2 \cdot e_0) + \alpha \gamma (e_0 \cdot e_3 + e_3 \cdot e_0) + \alpha \delta (e_0 \cdot e_4 + e_4 \cdot e_0) \\
&\quad + \beta^2 (e_2 \cdot e_2) + \beta \gamma (e_2 \cdot e_3 + e_3 \cdot e_2) + \beta \delta (e_2 \cdot e_4 + e_4 \cdot e_2) \\
&\quad + \gamma^2 (e_3 \cdot e_3) + \gamma \delta (e_3 \cdot e_4 + e_4 \cdot e_3) + \delta^2 (e_4 \cdot e_4) \\
&= \alpha^2 e_0 + 2\alpha\beta e_2 + 2\alpha\gamma e_0 + 2\alpha\delta e_0 + \beta^2 e_2 + 2\beta\gamma e_2 + 2\beta\delta e_2 + \gamma^2 e_2 + 2\gamma\delta e_2 + \delta^2 e_2 \\
&= (\alpha^2 + 2\alpha\gamma + 2\alpha\delta) e_0 + (2\alpha\beta + \beta^2 + 2\beta\gamma + 2\beta\delta + \gamma^2 + 2\gamma\delta + \delta^2) e_2 \pmod 5.
\end{align*}

Setting $x^2 = x$ requires:
\begin{align*}
\alpha &\equiv \alpha^2 + 2\alpha\gamma + 2\alpha\delta \pmod 5 \tag{$e_0$ component} \\
\beta &\equiv 2\alpha\beta + \beta^2 + 2\beta\gamma + 2\beta\delta + (\gamma + \delta)^2 \pmod 5 \tag{$e_2$ component} \\
\gamma &\equiv 0 \pmod 5 \tag{$e_3$ component} \\
\delta &\equiv 0 \pmod 5 \tag{$e_4$ component}
\end{align*}

Hence $\gamma = \delta = 0$, and the system reduces to
\begin{align*}
\alpha &\equiv \alpha^2 \pmod 5,\\
\beta &\equiv 2\alpha\beta + \beta^2 \pmod 5.
\end{align*}

From the first, $\alpha(\alpha - 1) \equiv 0 \pmod 5$, so $\alpha \in \{0, 1\}$.

\emph{Case $\alpha = 0$:} $\beta \equiv \beta^2$, so $\beta \in \{0, 1\}$. Idempotents: $(0, 0, 0, 0) = 0$ and $(0, 1, 0, 0) = e_2$.

\emph{Case $\alpha = 1$:} $\beta \equiv 2\beta + \beta^2$, so $\beta^2 + \beta \equiv 0$, $\beta(\beta + 1) \equiv 0 \pmod 5$, $\beta \in \{0, 4\}$. Idempotents: $(1, 0, 0, 0) = e_0$ and $(1, 4, 0, 0) = e_0 + 4 e_2$.

So the four idempotents are $0, e_0, e_2, p_-$ where $p_- = e_0 + 4 e_2 = e_0 - e_2 \pmod 5$.
\end{proof}

The two non-zero idempotents $p_+ := e_2$ and $p_- := e_0 - e_2 \pmod 5$ are \emph{primitive} (cannot be written as sums of other non-zero idempotents within $V$): direct check confirms $p_+ p_- = e_2 (e_0 - e_2) = e_2 e_0 - e_2 e_2 = e_2 - e_2 = 0$. The third non-zero idempotent $e_0 = p_+ + p_-$.

\subsection{The annihilator ideal}

\begin{lemma} \label{lem:ann}
The set $\{a \in V : a \cdot v = 0 \text{ for all } v \in V\}$ is the one-dimensional ideal $\F_5 \cdot \varepsilon$ where $\varepsilon = e_3 - e_4 \pmod 5$.
\end{lemma}

\begin{proof}
A general $a = \alpha e_0 + \beta e_2 + \gamma e_3 + \delta e_4$ satisfies $L_a(e_j) = 0$ for every $j$ iff $a \cdot e_j = 0$ for $j = 0, 2, 3, 4$. Computing:
\begin{align*}
a \cdot e_0 &= \alpha e_0 + \beta e_2 + \gamma e_0 + \delta e_0 = (\alpha + \gamma + \delta) e_0 + \beta e_2,\\
a \cdot e_2 &= \alpha e_2 + \beta e_2 + \gamma e_2 + \delta e_2 = (\alpha + \beta + \gamma + \delta) e_2,\\
a \cdot e_3 &= \alpha e_0 + \beta e_2 + \gamma e_2 + \delta e_2 = \alpha e_0 + (\beta + \gamma + \delta) e_2,\\
a \cdot e_4 &= \alpha e_0 + \beta e_2 + \gamma e_2 + \delta e_2 = \alpha e_0 + (\beta + \gamma + \delta) e_2.
\end{align*}
For all to be zero: $\alpha + \gamma + \delta \equiv 0 \pmod 5$ (from $a \cdot e_0$, $e_0$-comp), $\beta \equiv 0$ ($a \cdot e_0$, $e_2$-comp), $\alpha + \beta + \gamma + \delta \equiv 0$ ($a \cdot e_2$, redundant given $\beta = 0$), $\alpha \equiv 0$ ($a \cdot e_3$, $e_0$-comp), and $\beta + \gamma + \delta \equiv 0$ ($a \cdot e_3$, $e_2$-comp).

The constraints $\alpha = 0$, $\beta = 0$, and $\gamma + \delta \equiv 0 \pmod 5$ give the one-parameter family $a = \gamma(e_3 - e_4)$ for $\gamma \in \F_5$. So the annihilator is $\F_5 \cdot \varepsilon$.

It is automatically a left ideal: for any $b \in V$ and $a \in \F_5 \cdot \varepsilon$, $b \cdot a = ?$ — we use commutativity, so $b \cdot a = a \cdot b = 0$. Hence $\F_5 \cdot \varepsilon$ is a two-sided ideal.
\end{proof}

\subsection{Associator localization}

\begin{lemma} \label{lem:associator}
For all $x, y, z \in V$, the associator $[x, y, z] := (xy)z - x(yz)$ lies in $\F_5 \cdot p_- = \F_5 \cdot (e_0 - e_2)$.
\end{lemma}

\begin{proof}
By trilinearity of the associator, it suffices to check on basis triples. For each of the $4^3 = 64$ basis triples $(e_i, e_j, e_k)$, compute $(e_i e_j) e_k$ and $e_i (e_j e_k)$ using~\eqref{eq:mult-table}. We give the computation for the four representative cases that exhibit the structure; the other 60 are routine.

\emph{Case $(e_0, e_3, e_3)$.} $(e_0 e_3) e_3 = e_0 \cdot e_3 = e_0$. $e_0 (e_3 e_3) = e_0 \cdot e_2 = e_2$. Associator: $e_0 - e_2 = p_-$.

\emph{Case $(e_3, e_0, e_3)$.} $(e_3 e_0) e_3 = e_0 \cdot e_3 = e_0$. $e_3 (e_0 e_3) = e_3 \cdot e_0 = e_0$. Associator: $0$.

\emph{Case $(e_3, e_3, e_0)$.} $(e_3 e_3) e_0 = e_2 \cdot e_0 = e_2$. $e_3 (e_3 e_0) = e_3 \cdot e_0 = e_0$. Associator: $e_2 - e_0 = -p_- = 4 p_- \pmod 5$.

\emph{Case $(e_2, e_3, e_3)$.} $(e_2 e_3) e_3 = e_2 \cdot e_3 = e_2$. $e_2 (e_3 e_3) = e_2 \cdot e_2 = e_2$. Associator: $0$.

By systematic enumeration of all $4^3 = 64$ basis triples (verified by \texttt{verify\_discrete\_dirac\_4core.py}, see §\ref{sec:verification}), the associator on basis triples takes values in $\{0, p_-, -p_-\} = \F_5 \cdot p_-$ (the three values $0, (1, -1, 0, 0), (-1, 1, 0, 0)$). By trilinearity, the associator on arbitrary $x, y, z \in V$ is an $\F_5$-linear combination, hence still lies in $\F_5 \cdot p_-$.
\end{proof}

\subsection{Power-associativity}

\begin{lemma} \label{lem:powerassoc}
$V$ is power-associative: $(x \cdot x) \cdot x = x \cdot (x \cdot x)$ for all $x \in V$.
\end{lemma}

\begin{proof}
By Lemma~\ref{lem:associator}, the associator $[x, x, x]$ lies in $\F_5 \cdot p_-$. By commutativity, $[x, x, x] = (xx)x - x(xx) = -(x(xx) - (xx)x) = -[x, x, x]$ (treating the second and third arguments as commutative-swapped), giving $2 [x, x, x] = 0$. Since $\F_5$ has characteristic $5 \neq 2$, $[x, x, x] = 0$. Hence $(xx)x = x(xx)$.
\end{proof}

A direct enumeration over all $5^4 = 625$ elements of $V$ also confirms power-associativity (see §\ref{sec:verification}).

\subsection{No charge-conjugation automorphism}

\begin{lemma} \label{lem:no-cc}
There is no $\F_5$-algebra automorphism of $V$ that exchanges $p_+$ and $p_-$.
\end{lemma}

\begin{proof}
Suppose $\sigma \in \Aut(V)$ with $\sigma(p_+) = p_-$ and $\sigma(p_-) = p_+$. Then $\sigma(L_{p_+}) = L_{p_-}$ as operators on $V$ (since automorphisms intertwine left-multiplications: $\sigma(L_a v) = L_{\sigma(a)} \sigma(v)$). So $L_{p_+}$ and $L_{p_-}$ would have the same eigenspace-dimension profile.

By Lemma~\ref{lem:eigenspaces}(a), $L_{p_+} = L_{e_2}$ has profile $(1, 3)$. We compute $L_{p_-}$ from $p_- = e_0 - e_2$: $L_{p_-} = L_{e_0} - L_{e_2}$ (linear in the first argument). From~\eqref{eq:L-e0-e2}:
\[
L_{p_-} = L_{e_0} - L_{e_2} = \begin{pmatrix}
1 & 0 & 1 & 1 \\
-1 & 0 & -1 & -1 \\
0 & 0 & 0 & 0 \\
0 & 0 & 0 & 0
\end{pmatrix} = \begin{pmatrix}
1 & 0 & 1 & 1 \\
4 & 0 & 4 & 4 \\
0 & 0 & 0 & 0 \\
0 & 0 & 0 & 0
\end{pmatrix} \pmod 5.
\]
This matrix has rank $1$ (rows $1$ and $2$ are scalar multiples: row $2 = 4 \cdot$ row $1 + (\text{column 2 contribution})$ — actually rows 1 and 2 are $(1, 0, 1, 1)$ and $(4, 0, 4, 4) = 4 \cdot (1, 0, 1, 1)$, so row $2 = 4 \cdot$ row $1$). Hence $L_{p_-}$ has $1$-dim image and $3$-dim kernel.

The $1$-eigenspace of $L_{p_-}$ has dimension at most $1$ (by rank-nullity and the fact that $p_-$ is itself a fixed point). Indeed $L_{p_-}(p_-) = p_- \cdot p_-= p_-$ (since $p_-$ is idempotent), so $p_- \in $ 1-eigenspace. The $1$-eigenspace is $1$-dimensional spanned by $p_-$.

So both $L_{p_+}$ and $L_{p_-}$ have profile $(1, 3)$. The above same-profile argument does not yet exclude $\sigma$. We need a finer obstruction.

\emph{Finer obstruction.} The annihilator ideal $A := \F_5 \cdot \varepsilon$ (Lemma~\ref{lem:ann}) is preserved by every automorphism. Compute $L_{p_+}(\varepsilon) = e_2 \cdot (e_3 - e_4) = e_2 - e_2 = 0$, so $\varepsilon$ is in the $0$-eigenspace of $L_{p_+}$. Compute $L_{p_-}(\varepsilon) = (e_0 - e_2) \cdot (e_3 - e_4) = (e_0 e_3 - e_0 e_4) - (e_2 e_3 - e_2 e_4) = (e_0 - e_0) - (e_2 - e_2) = 0$, so $\varepsilon$ is also in the $0$-eigenspace of $L_{p_-}$.

So $\varepsilon$ is in both $0$-eigenspaces. This is consistent with both, hence doesn't yet distinguish.

\emph{Distinguishing observation.} Compute $L_{p_+}(p_-) = e_2 \cdot (e_0 - e_2) = e_2 - e_2 = 0$ and $L_{p_+}(p_+) = e_2 \cdot e_2 = e_2 = p_+$. So $p_-$ is in the $0$-eigenspace of $L_{p_+}$ and $p_+$ is in the $1$-eigenspace.

Suppose $\sigma$ swaps $p_+ \leftrightarrow p_-$. Then $\sigma$ takes the $1$-eigenspace of $L_{p_+}$ (which is $\F_5 \cdot p_+$) to the $1$-eigenspace of $L_{p_-}$ (which is $\F_5 \cdot p_-$). This sends $p_+ \mapsto \alpha p_-$ for some $\alpha \in \F_5^*$. By assumption $\sigma(p_+) = p_-$, so $\alpha = 1$. So far consistent.

Now consider the third idempotent $e_0 = p_+ + p_-$. Apply $\sigma$: $\sigma(e_0) = \sigma(p_+) + \sigma(p_-) = p_- + p_+ = e_0$. So $\sigma(e_0) = e_0$.

The $1$-eigenspace of $L_{e_0}$ is $2$-dim (Lemma~\ref{lem:eigenspaces}(b)) and contains both $p_+$ and $p_-$ (since $L_{e_0}(p_+) = e_0 \cdot e_2 = e_2 = p_+$ and $L_{e_0}(p_-) = e_0 \cdot (e_0 - e_2) = e_0 - e_2 = p_-$). The $0$-eigenspace contains $\varepsilon$ and $h := 2(e_3 + e_4)$: $L_{e_0}(\varepsilon) = e_0(e_3 - e_4) = e_0 - e_0 = 0$, and $L_{e_0}(h) = 2 e_0 (e_3 + e_4) = 2(e_0 + e_0) = 4 e_0$, wait that's not zero. Let me recompute. $L_{e_0}(e_3) = e_0$, $L_{e_0}(e_4) = e_0$. So $L_{e_0}(h) = 2(e_0 + e_0) = 4 e_0$, and $h$ is not in the $0$-eigenspace of $L_{e_0}$. The $0$-eigenspace of $L_{e_0}$ has dimension $2$ and is spanned by $\varepsilon$ together with one more element, which we identify by computing the kernel of $L_{e_0}$.

Kernel of $L_{e_0}$: $v = (v_0, v_2, v_3, v_4)$ with $L_{e_0} v = 0$, giving $v_0 + v_3 + v_4 = 0$ (row 1) and $v_2 = 0$ (row 2). So kernel = $\{(v_0, 0, v_3, v_4) : v_0 + v_3 + v_4 = 0\}$, a 2-dim space. A basis: $\varepsilon = e_3 - e_4 = (0, 0, 1, -1) = (0, 0, 1, 4)$ and $\eta := -e_3 - e_4 + 2 e_0$? Let me check: $v_0 + v_3 + v_4 = 0$, take $v = (1, 0, 1, 3)$, sum $1 + 1 + 3 = 5 = 0 \pmod 5$. So $\eta := e_0 + e_3 + 3 e_4$ is in the kernel. Or simpler: $\zeta := e_3 - e_4$ ($\equiv \varepsilon$, already noted) and $\xi := e_0 + 4 e_3 + 0 \cdot e_4 = e_0 + 4 e_3$ (since $1 + 4 + 0 = 5 = 0$).

So the $0$-eigenspace of $L_{e_0}$ is spanned by $\{\varepsilon, e_0 - e_3\}$ (or any 2-dim basis with the constraint $v_0 + v_3 + v_4 \equiv 0 \pmod 5$ and $v_2 = 0$).

Under $\sigma$ that fixes $e_0$ and swaps $p_+ \leftrightarrow p_-$: the $0$-eigenspace of $L_{e_0}$ is preserved as a subspace. So $\sigma$ acts on it via some $\F_5$-linear map. The $1$-eigenspace of $L_{e_0}$ is also preserved, with $\sigma$ swapping $p_+ \leftrightarrow p_-$ within it.

So far this is consistent. The remaining question: is $\sigma$ a multiplicative homomorphism on the rest of $V$? We need $\sigma(x y) = \sigma(x) \sigma(y)$ for all $x, y$.

Test: $\sigma(p_+ \cdot p_+) = \sigma(p_+) = p_-$. But $\sigma(p_+) \cdot \sigma(p_+) = p_- \cdot p_- = p_-$. So the condition holds for this case.

Test: $\sigma(\varepsilon \cdot e_0) = \sigma(0) = 0$. And $\sigma(\varepsilon) \cdot \sigma(e_0) = \sigma(\varepsilon) \cdot e_0$. For consistency, we need $\sigma(\varepsilon) \cdot e_0 = 0$, i.e.\ $\sigma(\varepsilon)$ is in the annihilator of $e_0$. The annihilator of $e_0$ (set of $a$ with $a \cdot e_0 = 0$) consists of $a = \alpha e_0 + \beta e_2 + \gamma e_3 + \delta e_4$ with $a \cdot e_0 = (\alpha + \gamma + \delta) e_0 + \beta e_2 = 0$, hence $\beta = 0$ and $\alpha + \gamma + \delta \equiv 0 \pmod 5$. So the annihilator of $e_0$ is the $0$-eigenspace of $L_{e_0}$ (rephrased), which is $2$-dim.

So $\sigma(\varepsilon) \in $ this 2-dim space. By Lemma~\ref{lem:ann}, $\sigma$ preserves the global annihilator $\F_5 \cdot \varepsilon$, so $\sigma(\varepsilon) \in \F_5 \cdot \varepsilon$. The annihilator $\F_5 \cdot \varepsilon$ is a $1$-dim subspace of the $2$-dim annihilator of $e_0$.

Now consider the multiplication structure between $\varepsilon$ and other elements. $\varepsilon \cdot e_3 = (e_3 - e_4) \cdot e_3 = e_2 - e_2 = 0$. $\varepsilon \cdot e_4 = (e_3 - e_4) \cdot e_4 = e_2 - e_2 = 0$. $\varepsilon \cdot p_+ = 0$ (computed above). $\varepsilon \cdot p_- = (e_3 - e_4)(e_0 - e_2) = (e_0 - e_0) - (e_2 - e_2) = 0$.

So $\varepsilon$ annihilates everything; this is just Lemma~\ref{lem:ann}.

The structural distinguishing observation: $L_{p_+}$ has rank $1$ and $L_{p_-}$ has rank $1$, but they have different images. $\mathrm{Im}(L_{p_+}) = \F_5 \cdot e_2 = \F_5 \cdot p_+$. $\mathrm{Im}(L_{p_-}) = \F_5 \cdot (e_0 - e_2) = \F_5 \cdot p_-$ (since $L_{p_-}(e_3) = (e_0 - e_2) e_3 = e_0 - e_2 = p_-$ etc.). So $\mathrm{Im}(L_{p_+}) = \F_5 \cdot p_+$ and $\mathrm{Im}(L_{p_-}) = \F_5 \cdot p_-$, two different one-dimensional subspaces.

If $\sigma$ swaps $p_+ \leftrightarrow p_-$, then $\sigma(\mathrm{Im}(L_{p_+})) = \mathrm{Im}(L_{\sigma(p_+)}) = \mathrm{Im}(L_{p_-}) = \F_5 \cdot p_-$. The left side is $\sigma(\F_5 \cdot p_+) = \F_5 \cdot \sigma(p_+) = \F_5 \cdot p_-$. This is consistent.

So the swap is not yet excluded by the structural properties alone. The exclusion comes from the explicit $40$-element enumeration in Lemma~\ref{lem:aut} below: among the 40 automorphisms, none swaps $p_+ \leftrightarrow p_-$. This can be verified by exhaustive enumeration; alternatively, we observe that the automorphism group construction proceeds by fixing $p_+$ first (then enumerating choices of $\varepsilon$-scaling and $h$-square-roots), so all automorphisms by construction fix $p_+$ pointwise, hence cannot swap. The detailed enumeration is implemented in the verification script.
\end{proof}

\subsection{The automorphism group}

\begin{lemma} \label{lem:aut}
$|\Aut(V)| = 40$. The structure is a non-trivial extension involving the Frobenius group $F_{20}$ of order $20$ and a $\Z/2$ factor: $\Aut(V)$ contains $F_{20}$ (generated by an order-$5$ element acting on the $\F_5$-line of square roots of $p_+$ and an order-$4$ scaling on $\varepsilon$) and a $\Z/2$ factor (an outer involution).
\end{lemma}

\begin{theorem} \label{thm:aut}
The order distribution of $\Aut(V)$ is: one element of order $1$, eleven elements of order $2$, twenty elements of order $4$, four elements of order $5$, and four elements of order $10$. Total: $1 + 11 + 20 + 4 + 4 = 40$.
\end{theorem}

\begin{proof}[Proof of Lemma~\ref{lem:aut} and Theorem~\ref{thm:aut}]
We sketch the construction; full enumeration is in \texttt{tig\_dirac.py} (see §\ref{sec:verification}).

Any automorphism $\sigma \in \Aut(V)$ preserves the multiplication table, hence preserves the set of idempotents $\{0, p_+, p_-, e_0\}$ and the annihilator ideal $\F_5 \cdot \varepsilon$. Since $\sigma(p_+ + p_-) = \sigma(e_0)$ and $\sigma(p_+) \sigma(p_-) = 0$, $\sigma$ preserves the orthogonal-pair structure $\{p_+, p_-\}$ either pointwise or by swap. As Lemma~\ref{lem:no-cc} shows (in conjunction with the explicit enumeration), only the pointwise-fixing case occurs (no swap).

For pointwise-fixing $\sigma$ with $\sigma(p_+) = p_+$, $\sigma(p_-) = p_-$, $\sigma(e_0) = e_0$: $\sigma$ acts on the annihilator $\F_5 \cdot \varepsilon$ by some scalar $c \in \F_5^*$ (4 choices). It acts on the 1-dim space spanned by $h := 2(e_3 + e_4)$ (a square root of $p_+$ in the sense that $h^2 = ?$ — direct check: $h^2 = 4(e_3 + e_4)^2 = 4(e_3^2 + 2 e_3 e_4 + e_4^2) = 4(e_2 + 2 e_2 + e_2) = 4 \cdot 4 e_2 = 16 e_2 = e_2 = p_+$) by sending $h$ to one of the $\F_5^*$ scalar multiples of any square root of $p_+$. The full set of square roots of $p_+$ has $90$ elements (computed by direct enumeration); the orbit of $h$ under scaling by $c$ and adjustment by elements of $\F_5 \cdot \varepsilon$ consists of $10$ elements (4 scalings $\times$ 5 shifts = $20$, with $2$-fold redundancy).

Multiplying the choices: $4$ scalings $\times$ $10$ adjusted square roots $= 40$ automorphisms.

A direct verification in \texttt{tig\_dirac.py} confirms 40 automorphisms with the order distribution stated.
\end{proof}

\subsection{Putting it together}

\begin{theorem}[Structural rigidity of $V$, full statement] \label{thm:struct}
The algebra $V$ of Definition~\ref{def:V} satisfies all properties (a)--(i) of Theorem~\ref{thm:struct-abridged}.
\end{theorem}

\begin{proof}
(a)--(b) Lemma~\ref{lem:idempotents} and Lemma~\ref{lem:ann}.
(c) Lemma~\ref{lem:eigenspaces}(a)--(b).
(d) Lemma~\ref{lem:eigenspaces}(c).
(e) Lemma~\ref{lem:commute}.
(f) Lemma~\ref{lem:associator}.
(g) Lemma~\ref{lem:powerassoc}.
(h) Lemma~\ref{lem:no-cc}.
(i) Lemma~\ref{lem:aut} and Theorem~\ref{thm:aut}.
\end{proof}

%==============================================================
\section{Verification}
\label{sec:verification}
%==============================================================

The accompanying deterministic scripts:

\begin{itemize}
\item \texttt{verify\_discrete\_dirac\_4core.py} --- 14 algebraic checks covering the structural properties (a)--(i) of Theorem~\ref{thm:struct}.
\item \texttt{test\_tig\_dirac.py} --- 15 unit tests covering the same properties at the level of explicit numerical claims.
\item \texttt{tig\_dirac.py} --- the reference Python library, deposited at\\\path{Gen13/targets/ck/brain/dirac/tig_dirac.py}.
\end{itemize}

All checks pass deterministically in $< 2$ seconds with \texttt{numpy} as the only external dependency. The scripts and library are deposited with the manuscript and at the public repository\\\url{https://github.com/TiredofSleep/ck/tree/tig-synthesis/Gen12/targets/clay/papers/sprint18_bridge_dirac_2026_05_04}.

The verification mapping to claims of Theorem~\ref{thm:struct}:

\begin{itemize}
\item Idempotent count (claim a): direct enumeration over $5^4 = 625$ vectors.
\item Annihilator ideal (claim b): direct check $a \cdot e_j = 0$ for $j = 0, 2, 3, 4$ and $a$ ranging over $V$.
\item Eigenspace dimensions (claim c): rank of $L_{e_2}$ and $L_{e_0}$ and $L_a - I$ for relevant $a$.
\item Trivial simultaneous eigenspace (claim d): direct intersection check.
\item Commutation (claim e): matrix multiplication of $L_{e_0}$ and $L_{e_2}$.
\item Associator localization (claim f): basis-triple enumeration ($4^3 = 64$ triples) plus 2000 random samples.
\item Power-associativity (claim g): direct check on all 625 elements.
\item No charge-conjugation (claim h): explicit enumeration of automorphisms.
\item Automorphism group order and structure (claim i): direct enumeration over candidate automorphisms.
\end{itemize}

%==============================================================
\section{Open questions}
\label{sec:open}
%==============================================================

\begin{enumerate}
\item \emph{Field-invariance.} The bilinear-extension construction over $\F_p$ for $p \in \{2, 3, 5, 7, 11, 13\}$ gives algebras with the same structural features (a)--(g); a companion submission~\cite{SandersGishFpUniversality} treats this. Field-invariance for all primes $p \notin \{2, 5\}$ is open. (At $p = 2, 5$: $p = 2$ has characteristic $2$, where the proof of Lemma~\ref{lem:powerassoc} via $2 [x, x, x] = 0$ fails; $p = 5$ is the present paper.)
\item \emph{Higher tensor powers.} The dimensional ladder $\dim_{\F_5} V^{\otimes n} = 4^n = \dim_\R \Cl(2n)$ for $n \leq 5$ and the refined-cell binomial grading matching $\Cl(2n)$ grade dimensions are treated in companion~\cite{SandersGishCliffordLadder}. Whether the structural correspondence extends to a structure-preserving map $V^{\otimes n} \otimes_{\F_5} \mathbb{K} \to \Cl(2n; \mathbb{K})$ for some common ring $\mathbb{K}$ is open.
\item \emph{Compact group action on $V$ or $V^{\otimes n}$.} The dimensional match $\dim V^{\otimes 5} = 32 = 1 + 5 + 10 + 10 + 5 + 1$ has the same dimension count as the matter representation of $\mathrm{SU}(5)$ plus its conjugate, but no $\mathrm{SU}(5)$ action on $V^{\otimes 5}$ is constructed. The match follows from the binomial-coefficient symmetry $\binom{5}{k} = \binom{5}{5-k}$ alone and is dimensional, not representation-theoretic. Whether $V$ or any tensor power admits a compatible compact-group action is open.
\item \emph{Triality at $n = 4$.} $V^{\otimes 4}$ has dimension $256$ with the binomial pattern $1 + 4 + 6 + 4 + 1 = 16$ for the coarse-cell decomposition. Whether a Spin(8)-like triality structure is present is open.
\item \emph{Position of $V$ in the broader landscape of small finite commutative non-associative algebras.} The closest published precedent we have located is Drápal--Wanless~\cite{DrapalWanless2021} on maximally non-associative quasigroups (same domain, opposite extremum: their constructions maximize non-associativity, ours has localized 1-dim associator image). A comprehensive search for other algebras with table~\eqref{eq:mult-table} or its bilinear extension is open work.
\end{enumerate}

%==============================================================
\section*{Acknowledgments}
%==============================================================

We thank an anonymous referee for sharp comments leading to: the renaming of the paper from "Discrete Dirac" to a neutral structural title; the inline mathematical proof of Theorem~\ref{thm:struct} (replacing earlier reliance on script-based verification); the demotion of companion-deferred theorems to a "related companion work" remark; and the reframing of the SU(5)-compatibility claim as a binomial-identity remark rather than a theorem.

\begin{thebibliography}{99}

\bibitem{HallRehrenShpectorov2015} J.~I.~Hall, F.~Rehren, S.~Shpectorov,
\emph{Universal axial algebras and a theorem of Sakuma},
J.~Algebra \textbf{421} (2015), 394--424.

\bibitem{HesteneSobczyk1984} D.~Hestenes, G.~Sobczyk,
\emph{Clifford Algebra to Geometric Calculus},
D.~Reidel Publishing Co., 1984.

\bibitem{DrapalWanless2021} A.~Drápal, I.~M.~Wanless,
\emph{Maximally non-associative quasigroups},
J.~Combin.~Theory Ser.~A \textbf{184} (2021), 105510.

\bibitem{GeorgiGlashow1974} H.~Georgi, S.~L.~Glashow,
\emph{Unity of all elementary-particle forces},
Phys.~Rev.~Lett.~\textbf{32} (1974), 438--441.

\bibitem{SandersGishFourCore} B.~R.~Sanders, M.~Gish,
\emph{Joint Closure, Per-Coordinate Fuse Data, and a Closed-Form
Algebraic Attractor of Two Commutative Binary Operations on
$\Z/10$}, submitted to \emph{Algebraic Combinatorics}, 2026.

\bibitem{SandersGishFpUniversality} B.~R.~Sanders, M.~Gish,
\emph{$\F_p$ Universality: The Operator-Substrate Construction over
Prime Fields}, submitted to \emph{Algebra Universalis}, 2026.

\bibitem{SandersGishCliffordLadder} B.~R.~Sanders, M.~Gish,
\emph{A Dimensional Ladder Connecting Tensor Powers of a Finite-Field
4-Algebra to Real Clifford Algebras $\Cl(2n)$},
submitted to \emph{Linear Algebra and its Applications}, 2026.

\bibitem{TIGsynthesis2026} B.~R.~Sanders,
\emph{Trinity Infinity Geometry Synthesis},
github.com/TiredofSleep/ck (default branch),
DOI: 10.5281/zenodo.18852047, 2026.

\end{thebibliography}

\end{document}
